\par\Needspace{6\baselineskip}\medskip
\hypertarget{native-113-article084}{}\pdfbookmark[2]{84.}{article-084}
\begin{center}84.\end{center}
Quum per duos radios vectores magnitudine et positione datos atque elementum orbitae unum orbitam integram determinare liceat, per illa data etiam \SourceLeaf{114}\emph{tempus,} intra quod corpus coeleste ab uno radio vectore ad alterum movetur, determinabile erit, siquidem corporis massam vel negligimus vel saltem tamquam cognitam spectamus: nos suppositioni priori inhaerebimus, ad quam posterior facile reducitur. Hinc vice versa patet, duos radios vectores magnitudine et positione datos una cum tempore, intra quod corpus coeleste spatium intermedium describit, orbitam integram determinare. Hoc vero problema, ad gravissima in theoria motus corporum coelestium referendum, haud ita facile solvitur, quum expressio temporis per elementa transscendens sit, insuperque satis complicata. Eo magis dignum est, quod omni cura tractetur: quamobrem lectoribus haud ingratum fore speramus, quod praeter solutionem post tradendam, quae nihil amplius desiderandum relinquere videtur, eam quoque oblivioni eripiendam esse censuimus, qua olim antequam ista se obtulisset frequenter usi sumus. Problemata difficiliora semper iuvat pluribus viis aggredi, nec bonam spernere etiamsi meliorem praeferas. Ab expositione huius methodi anterioris initium facimus.
